\documentclass[11pt,oneside]{article}
\usepackage[a4paper,margin=27mm,headheight=15pt]{geometry}
\usepackage[T1]{fontenc}
\usepackage{lmodern,microtype}
\usepackage{amsmath,amssymb,amsthm,mathtools,bm,mathrsfs}
\usepackage{booktabs,longtable,array,enumitem,xurl,fancyhdr}
\usepackage[hidelinks]{hyperref}
\usepackage{bookmark}
\setlength{\parskip}{3pt}
\setlength{\parindent}{1.2em}
\setlength{\emergencystretch}{3em}
\numberwithin{equation}{section}
\newtheorem{theorem}{Theorem}[section]
\newtheorem{proposition}[theorem]{Proposition}
\newtheorem{lemma}[theorem]{Lemma}
\newtheorem{corollary}[theorem]{Corollary}
\theoremstyle{definition}
\newtheorem{definition}[theorem]{Definition}
\newtheorem{example}[theorem]{Example}
\theoremstyle{remark}
\newtheorem{remark}[theorem]{Remark}
\DeclareMathOperator{\Li}{Li}
\DeclareMathOperator{\CT}{CT}
\DeclareMathOperator{\FP}{FP}
\DeclareMathOperator{\Rea}{Re}
\DeclareMathOperator{\cont}{cont}
\newcommand{\dd}{\,\mathrm d}
\newcommand{\N}{\mathbb N}
\newcommand{\Q}{\mathbb Q}
\newcommand{\C}{\mathbb C}
\newcommand{\R}{\mathbb R}
\newcommand{\uu}{\mathbf u}
\newcommand{\mm}{\mathbf m}
\newcommand{\ee}{\mathbf e}
\newcommand{\CC}{\mathcal C}
\newcommand{\KK}{\mathcal K}
\newcommand{\Acal}{\mathcal A}
\newcommand{\RR}{\mathcal R}
\newcommand{\Gop}{\mathcal G}
\newcommand{\Top}{\mathcal T}
\newcommand{\spec}[2]{\zeta^{(#1)}(#2,a)}
\newcommand{\src}[1]{\nolinkurl{#1}}
\pagestyle{fancy}
\fancyhf{}
\fancyhead[L]{\small Endpoint regularization of harmonic polylogarithms}
\fancyhead[R]{\small ProveIt continuation}
\fancyfoot[C]{\thepage}
\renewcommand{\headrulewidth}{0.3pt}
\setcounter{tocdepth}{2}
\hypersetup{pdftitle={Endpoint Regularization of Logarithmic Harmonic Polylogarithms},pdfauthor={Research continuation prepared with OpenAI assistance}}

\begin{document}
\begin{titlepage}
\thispagestyle{empty}
\vspace*{16mm}
\noindent\textsc{ProveIt / Analysis / Polylogarithms}\par
\vspace{14mm}
\noindent{\LARGE\bfseries Endpoint Regularization of\par
Logarithmic Harmonic Polylogarithms\par}
\vspace{7mm}
\noindent{\Large Stieltjes coefficients, Gamma transfer,\par
spectral finite parts, and exact antiderivatives\par}
\vspace{14mm}
\noindent Research continuation prepared for Vladimir Reshetnikov\par
\noindent with OpenAI assistance\par
\noindent 10 October 2026\par
\vspace{12mm}
\begin{abstract}
We give an explicit coefficient calculus for the endpoint constants of
symmetrized, logarithmically decorated harmonic polylogarithms on a common
Hurwitz lattice. A normally convergent canonical product expresses the
sharp-cutoff polynomials through generalized Stieltjes constants and spectral
derivatives of Hurwitz zeta. A proved Laplace--Stieltjes transfer converts
these polynomials to Abel endpoint polynomials by the finite operator
$\Gamma(1+\partial_L)$. This yields formulas at every depth and every fixed
logarithmic decoration, including a single-decoration family in which all
positive-index Stieltjes constants cancel at the unshifted lattice. Explicit
examples involve $\zeta'(2)$, $\zeta''(2)$, $\zeta'(3)$ and ordinary zeta products.
We separately compute diagonal spectral finite parts and their ray-dependent
correction terms: they are not the same regularization. A beta-weighted moment law evaluates the full unweighted family at every
depth. Convergent repeated antiderivatives admit hypergeometric generators; at every positive integer
shift a finite deletion formula reduces the entire unweighted hierarchy to
the classical height-one Gamma quotient. Rational shifts give cyclotomic
identities with an explicit logarithmic scale correction.

The symmetric-sum principle, Gamma regularization, and the height-one Gamma
identity are classical ingredients, not claimed discoveries. The contribution
is the unified, normalization-explicit shifted and decorated calculus, its
consequences, and its reproducible implementation. All stated universal
identities have proofs; finite symbolic tests and numerical diagnostics are
kept distinct from those proofs. No arithmetic independence, resolution of
the repository's $S_6$ or $S_8$ conjectures, or proof-assistant formalization
is claimed.
\end{abstract}
\vfill
\noindent\small Repository snapshot:\par
\noindent\src{fc4d3bf80534ad7c901d3b8c9e71baf2df0064ed}\par
\noindent Complete source, validation scripts, recorded outputs and audit notes
accompany this article.
\end{titlepage}

\begingroup
\small
\setlength{\parskip}{0pt}
\tableofcontents
\endgroup
\newpage

\section{Research question, scope, and relation to existing work}

The polylogarithm programme in ProveIt connects harmonic sums with Gamma
functions, Hurwitz-zeta jets, generalized Stieltjes constants and integration
operators. Its canonical manuscript already distinguishes proved reductions
from numerical candidates and, in particular, retains the conjectural status
of $S_6$ and $S_8$~\cite{repo}. We continue a different exact-identity question:

\begin{quote}
For a fully symmetrized harmonic sum with prescribed powers of logarithms,
can one compute the entire endpoint subtraction polynomial, and hence its
constant term, at arbitrary depth and a common positive Hurwitz shift?
How does the answer change when the regularization is changed?
\end{quote}

The answer is affirmative for the family defined below. The coefficients
require only finitely many single-variable Hurwitz-zeta derivatives and
Stieltjes constants. This is a statement about the explicitly symmetrized
family, not a reduction theorem for every ordered multiple polylogarithm.
The distinction is essential: symmetrization is the step that removes the
unrestricted nested-order information.

There are three layers of results. First, Theorems~\ref{thm:cutoff}
and~\ref{thm:abel} give all-depth cutoff and Abel polynomials. Second,
Theorem~\ref{thm:single-decoration} gives an exact generating function for one
arbitrary logarithmic decoration and proves cancellation of positive-index
Stieltjes constants at $a=1$ in every depth greater than one. Third,
Theorems~\ref{thm:diagonal}, \ref{thm:beta-moments} and~\ref{thm:integer-primitives}
give explicit regularization anomalies, all-depth beta moments and a
convergent antiderivative hierarchy.
The remaining results supply higher decorated examples, shift laws,
differential recurrences and rational-lattice transport.

\subsection{What is classical and what is developed here}

The finite symmetric-sum identity is an elementary form of the symmetric
multiple-zeta relations associated with Hoffman~\cite{Hoffman}; it can be
proved without invoking multiple-zeta continuation. The comparison between
harmonic and shuffle regularization is developed by Ihara, Kaneko and
Zagier~\cite{IKZ}. Kaneko and Sakata give a particularly useful formulation in
terms of cutoff and polylogarithmic asymptotic polynomials~\cite{KS}. Our
cutoff variable is $\log N$, rather than $\log N+\gamma$. This choice accounts
for the Euler-constant translation in the Gamma operator below.

The height-one generating identity
\[
1-\frac{\Gamma(1-u)\Gamma(1-v)}{\Gamma(1-u-v)}
=\sum_{p,d\ge1}\zeta(p+1,\{1\}^{d-1})v^p u^d
\]
is classical, following Aomoto and Drinfeld; see~\cite{Aomoto,Drinfeld,KS}.
We give a beta-integral proof because its finite-shift extension requires
careful normalization. We do not present the unshifted identity as new.
Likewise, the Hurwitz-zeta and Gamma identities used in the proofs are
standard~\cite{DLMF}.

The present continuation develops these ingredients simultaneously for
arbitrary finite logarithmic content and a common positive shift, states the
complete subtraction polynomials, compares them with diagonal and general-ray
spectral constants, and makes the single-decoration cancellation explicit.
These formulas are proved here, but no claim of worldwide priority is made.
The literature and repository review was targeted rather than exhaustive.

\subsection{Repository review boundary}

The inspected material consists of the canonical manuscript's README and
main driver; selected portions of its harmonic-depth and integration
chapters; the log-Gamma-moment development; the alternating-harmonic
continuation README; and the incoming directory listing and intake
instructions. The alternating-harmonic continuation already treats
$[u^r]\prod_{j\le n}(1+u/j)$ through incomplete-beta reflection~\cite{repoalt}.
Our subject is nonalternating endpoint regularization with independent
logarithmic decorations, not a new proof of that reflection theorem.

The five incoming ZIP filenames were visible, but their binary payloads were
not inspected. Consequently this article does not certify nonoverlap with
those unexamined payloads, nor does it audit all 417 pages of the canonical
manuscript. The provenance and audit files specify this boundary. No original
repository file has been changed.

\section{Definitions and the one-variable input}
\label{sec:definitions}

Throughout, $a>0$ is real. Every logarithm of a lattice point $n+a$ is the
real logarithm. We use the descending-index convention
\begin{equation}\label{eq:mpl}
\Li_{s_1,\ldots,s_d}(z_1,\ldots,z_d)
=\sum_{k_1>\cdots>k_d\ge1}
 \frac{z_1^{k_1}\cdots z_d^{k_d}}{k_1^{s_1}\cdots k_d^{s_d}}.
\end{equation}
The references using ascending indices must therefore be read with their
indices reversed. When only $z_1=z$ is displayed, the other arguments are one.
For $|z|<1$, the exponential damping in the largest index justifies every
fixed spectral differentiation used here, locally uniformly in the orders.

Write
\begin{equation}\label{eq:stieltjes}
\zeta(1+t,a)=\frac1t+
 \sum_{r\ge0}\frac{(-1)^r\gamma_r(a)}{r!}t^r,
\qquad \gamma_0(a)=-\psi(a),\qquad \gamma=\gamma_0(1).
\end{equation}
A prime on $\zeta$ always differentiates its first, spectral argument;
$\partial_a$ will always be written explicitly. Define
\begin{equation}\label{eq:harmonic}
H_N^{(s)}(a)=\sum_{n=0}^{N-1}(n+a)^{-s},\qquad
h_{r,N}(a)=\sum_{n=0}^{N-1}\frac{\log^r(n+a)}{n+a}.
\end{equation}
Thus $h_{r,N}=(-1)^r\partial_s^rH_N^{(s)}(a)|_{s=1}$.
At $a=1$, $H_N^{(k)}(1)$ is the usual generalized harmonic number.
We use the rising Pochhammer symbol $(a)_n=\Gamma(a+n)/\Gamma(a)$,
with $(a)_0=1$.

\begin{lemma}[Exact finite sums and cutoff constants]\label{lem:harmonic}
For $r\ge0$ and $N\ge0$,
\begin{equation}\label{eq:hgamma}
h_{r,N}(a)=\gamma_r(a)-\gamma_r(N+a).
\end{equation}
As $N\to\infty$,
\begin{equation}\label{eq:h-asymptotic}
h_{r,N}(a)=\frac{\log^{r+1}(N+a)}{r+1}+\gamma_r(a)
 +O_a\!\left(\frac{(1+\log N)^r}{N}\right).
\end{equation}
For integers $k\ge2$ and $R\ge0$,
\begin{align}
\sum_{n=0}^{N-1}\frac{\log^R(n+a)}{(n+a)^k}
 &=(-1)^R\bigl(\zeta^{(R)}(k,a)-\zeta^{(R)}(k,N+a)\bigr),
 \label{eq:higher-power}\\
\sum_{n\ge0}\frac{\log^R(n+a)}{(n+a)^k}
 &=(-1)^R\zeta^{(R)}(k,a).\label{eq:higher-limit}
\end{align}
All estimates are uniform when $a$ stays in a compact subinterval of $(0,\infty)$.
\end{lemma}

\begin{proof}
Subtract the Laurent expansions in
$\zeta(s,a)-\zeta(s,N+a)=H_N^{(s)}(a)$ and compare coefficients at $s=1$.
The poles cancel and give~\eqref{eq:hgamma}, with the sign in
\eqref{eq:stieltjes} unchanged.

For the asymptotic, the first Euler summation formula gives, locally uniformly
for $t$ near zero,
\[
\zeta(1+t,x)=\frac{x^{-t}}{t}+\frac12x^{-1-t}
 +(1+t)\int_0^\infty \frac{\widetilde B_1(y)}{(x+y)^{2+t}}\dd y,
\]
where $\widetilde B_1(y)=1/2-\{y\}$; its values at integers are immaterial.
Only boundedness is used. Differentiation of the
integral $r$ times produces an error
$O(x^{-1}(1+\log x)^r)$; this follows by putting $y=xv$ and integrating
$(1+v)^{-2}$ times a fixed power of $\log(1+v)$. The constant term removed from
$x^{-t}/t$ has coefficient
$(-1)^{r+1}\log^{r+1}x/(r+1)!$. Comparison with~\eqref{eq:stieltjes} yields
\[
\gamma_r(x)=-\frac{\log^{r+1}x}{r+1}
 +O(x^{-1}(1+\log x)^r).
\]
Together with~\eqref{eq:hgamma}, this proves~\eqref{eq:h-asymptotic}.
The remaining identities follow from the absolutely convergent Hurwitz series
and spectral differentiation. The tail is
$O(N^{1-k}(1+\log N)^R)$ by an integral comparison. Uniformity in $a$ follows
from the same estimates.
\end{proof}

For integer $k\ge2$ the undecorated version is also a polygamma identity:
\begin{equation}\label{eq:polygamma-harmonic}
H_N^{(k)}(a)=\frac{(-1)^k}{(k-1)!}
 \bigl(\psi^{(k-1)}(a)-\psi^{(k-1)}(N+a)\bigr).
\end{equation}
This is one direct bridge between the generalized harmonic numbers and the
polygamma functions in the construction.

Fix an integer $R\ge0$. Put
\[
f_r(x)=\frac{\log^r x}{x},\quad
W_{\uu}(x)=\sum_{r=0}^R u_r f_r(x),\quad
V_{\uu}(L)=\sum_{r=0}^R\frac{u_r L^{r+1}}{r+1}.
\]
A \emph{content} is $\mm=(m_0,\ldots,m_R)\in\N_0^{R+1}$. Its depth,
logarithmic degree and total degree are
\begin{equation}\label{eq:degrees}
d=|\mm|=\sum_r m_r,\qquad
w(\mm)=\sum_r r m_r,\qquad
D(\mm)=d+w(\mm)=\sum_r(r+1)m_r.
\end{equation}
Let $\uu^{\mm}=\prod_r u_r^{m_r}$ and define
\begin{equation}\label{eq:finite-product}
e_{\mm,N}(a)=[\uu^{\mm}]
 \prod_{n=0}^{N-1}\bigl(1+W_{\uu}(n+a)\bigr).
\end{equation}
For $d>0$, this is a sum over strictly ordered lattice points and distinct
words of logarithmic powers with content $\mm$. Repeated decorations do
\emph{not} introduce an extra factorial. For example,
\[
e_{(1,1),N}(a)=\sum_{N>n>m\ge0}
 \frac{\log(n+a)+\log(m+a)}{(n+a)(m+a)},
\]
whereas $e_{(0,2),N}$ has numerator $\log(n+a)\log(m+a)$, not twice that
numerator.

For $0<z<1$ define the corresponding Abel sum
\begin{equation}\label{eq:E}
E_{\mm}(z;a)=\sum_{n_1>\cdots>n_d\ge0}z^{n_1+a}
 \sum_{\cont(r_1,\ldots,r_d)=\mm}\prod_{i=1}^d f_{r_i}(n_i+a).
\end{equation}
There is one summand for each distinct word of the indicated content. Define
$E_{\mathbf0}=1$ separately. At $a=1$,
\begin{equation}\label{eq:spectral-decorations}
E_{\mm}(z;1)=(-1)^{w(\mm)}
 \sum_{\cont(\mathbf r)=\mm}
 \left.\partial_{s_1}^{r_1}\cdots\partial_{s_d}^{r_d}
 \Li_{s_1,\ldots,s_d}(z,1,\ldots,1)\right|_{s_1=\cdots=s_d=1}.
\end{equation}
This places the entire coefficient calculus inside spectral derivatives of
multiple polylogarithms, rather than merely formal harmonic notation.

\section{The all-depth cutoff coefficient theorem}
\label{sec:cutoff}

Define a canonical product in the $R+1$ auxiliary variables:
\begin{equation}\label{eq:Cproduct}
\CC_a(\uu)=\exp\!\left(\sum_{r=0}^R\gamma_r(a)u_r\right)
 \prod_{n\ge0}\bigl(1+W_{\uu}(n+a)\bigr)
                    e^{-W_{\uu}(n+a)}.
\end{equation}
This is an ordinary entire function of finitely many complex variables. It
must not be confused with the formal Gamma-moment series introduced later.

\begin{theorem}[Cutoff polynomials and exact coefficient closure]
\label{thm:cutoff}
The product~\eqref{eq:Cproduct} converges normally on every compact subset of
$\C^{R+1}$. Near $\uu=0$, its logarithm is
\begin{align}
\log\CC_a(\uu)
={}&\sum_{r=0}^R\gamma_r(a)u_r\notag\\
&+\sum_{k\ge2}\frac{(-1)^{k-1}}{k}
 \sum_{r_1,\ldots,r_k=0}^R
 (-1)^{r_1+\cdots+r_k}\zeta^{(r_1+\cdots+r_k)}(k,a)
 u_{r_1}\cdots u_{r_k}.\label{eq:cumulants}
\end{align}
For every content $\mm$, set
\begin{equation}\label{eq:P}
P_{\mm,a}(L)=[\uu^{\mm}]\CC_a(\uu)e^{V_{\uu}(L)}.
\end{equation}
If $d>0$, this is a polynomial of degree exactly $D(\mm)$, with leading
coefficient
\[
\frac{1}{\prod_{r=0}^R m_r!(r+1)^{m_r}}.
\]
Moreover,
\begin{equation}\label{eq:cutoff-asymptotic}
e_{\mm,N}(a)=P_{\mm,a}(\log(N+a))
 +O_{a,\mm}\!\left(\frac{(1+\log N)^{D(\mm)}}{N}\right).
\end{equation}
In particular, the sharp-cutoff constant is
\begin{equation}\label{eq:Cconstant}
C_{\mm}(a):=P_{\mm,a}(0)=[\uu^{\mm}]\CC_a(\uu).
\end{equation}
Every coefficient of $P_{\mm,a}$ is a rational polynomial in the finitely many
atoms $\gamma_r(a)$ and $\zeta^{(j)}(k,a)$, with $k\le d$ and $j\le w(\mm)$;
only decorations present in $\mm$ supply Stieltjes atoms.
\end{theorem}

\begin{proof}
On a compact set of auxiliary variables,
$|W_{\uu}(n+a)|\le C(1+\log n)^R/n$ for all sufficiently large $n$.
Consequently $\sum_n|W_{\uu}(n+a)|^2$ converges uniformly there. For the tail
of the product, $\log(1+W)-W=O(|W|^2)$ uniformly; finitely many initial
factors are entire. This proves normal convergence and entire dependence,
including at zeros of individual factors. A sufficiently small polydisc
around zero avoids all those zeros, so expanding its logarithm is legitimate.
The power-sum identity~\eqref{eq:higher-limit} then proves
\eqref{eq:cumulants}.

For the asymptotic, take the logarithm of the \emph{finite} product as a formal
power series and retain total degree at most $d$:
\begin{equation}\label{eq:finite-log}
\log\prod_{n<N}(1+W_{\uu}(n+a))
=\sum_{k=1}^{d}\frac{(-1)^{k-1}}k
  \sum_{n<N}W_{\uu}(n+a)^k
  \pmod{\text{total degree}>d}.
\end{equation}
The $k=1$ coefficient is supplied by~\eqref{eq:h-asymptotic}; it equals
$V_{\uu}(\log(N+a))+\sum_r\gamma_r(a)u_r$ plus the stated errors. For $k\ge2$
the relevant finite power sum differs from its limit by
$O(N^{1-k}(1+\log N)^{r_1+\cdots+r_k})$.

Exponentiation and extraction of a fixed content use only finitely many
products. The other factors in such a product have at most the remaining
logarithmic degree. A singleton error contributes at most a constant times
$N^{-1}(1+\log N)^{D-1}$; a block of size $k\ge2$ contributes at most
$N^{1-k}(1+\log N)^{D-k}$. The weaker common error in
\eqref{eq:cutoff-asymptotic} is therefore valid. This reasoning is entirely
coefficientwise and does not assert uniform asymptotics for unrestricted
auxiliary variables.

The highest power of $L$ comes only from selecting all decorations in
$e^{V_{\uu}(L)}$, with the constant term of $\CC_a$. This gives the displayed
leading coefficient. The finite atom statement follows directly from
\eqref{eq:cumulants}. The empty content gives $P_{\mathbf0,a}=1$ exactly.
\end{proof}

\subsection{First exact cutoff formulas}

For compactness in examples, write $g_r=\gamma_r(a)$ and
$Z_k^{[j]}=\zeta^{(j)}(k,a)$. With only $u_0$ present,
\begin{align}
C_{(1)}(a)&=g_0,\notag\\
C_{(2)}(a)&=\frac{g_0^2-Z_2^{[0]}}2,\notag\\
C_{(3)}(a)&=\frac{g_0^3-3g_0Z_2^{[0]}+2Z_3^{[0]}}6.
\label{eq:C-first}
\end{align}
The mixed logarithmic coefficients begin with
\begin{align}
P_{(1,1),a}(L)
 &=(L+g_0)\left(\frac{L^2}{2}+g_1\right)+Z_2^{[1]},
 \label{eq:P11}\\
P_{(0,2),a}(L)
 &=\frac12\left(\frac{L^2}{2}+g_1\right)^2-\frac12 Z_2^{[2]}.
 \label{eq:P02}
\end{align}
For instance, the plus sign before $Z_2^{[1]}$ in~\eqref{eq:P11} comes from
subtracting the diagonal sum $\sum\log(n+a)/(n+a)^2=-\zeta'(2,a)$.
This sign is easy to lose when a logarithmic insertion is identified with an
unsigned spectral derivative.

If one assigns formal degree $r+1$ to $g_r$, degree $k+j$ to $Z_k^{[j]}$,
and degree one to $L$, every term in $P_{\mm,a}$ has degree $D(\mm)$.
This is useful bookkeeping for an implementation. It is not a statement
about a motivic weight filtration or arithmetic independence.

\subsection{Why symmetrization matters}

For labelled variables $x_{i,n}$, inclusion--exclusion on coincidences gives
\begin{equation}\label{eq:partition}
\sum_{\substack{n_1,\ldots,n_d\ge0\\n_i\ \mathrm{distinct}}}
 \prod_{i=1}^d x_{i,n_i}
=
\sum_{\pi\in\Pi_d}\mu(\pi)
 \prod_{B\in\pi}\left(\sum_{n\ge0}\prod_{i\in B}x_{i,n}\right),
\end{equation}
when the sums are finite or absolutely convergent. Here $\Pi_d$ is the set of
set partitions and
\[
\mu(\pi)=\prod_{B\in\pi}(-1)^{|B|-1}(|B|-1)!.
\]
One proof is to fix an equality pattern of the indices. Its coefficient on
the right is the sum of the partition-lattice M\"obius function below that
pattern, hence is one for the discrete pattern and zero otherwise. An
alternative elementary proof expands the finite logarithmic product in
\eqref{eq:finite-log} and uses the exponential formula for set partitions.

With $x_{i,n}=f_{r_i}(n+a)$, the labelled distinct-index sum is
$\prod_r m_r!$ times the content coefficient. This accounts for the absence
of any additional factorial in~\eqref{eq:finite-product}. An individual
ordered word has no such general reduction. We never infer its reduction
from that of the complete symmetrization.

\section{Gamma transfer from cutoff to Abel endpoints}
\label{sec:abel}

Set $\epsilon=-\log z>0$ and $L=\log(1/\epsilon)$. There is no ambiguity in
these real variables. The more customary $L_z=-\log(1-z)$ satisfies
$L_z=L+O(\epsilon)$, so replacing one by the other changes a fixed polynomial
by $o(1)$ as $z\uparrow1$.

\begin{lemma}[Polynomial Laplace--Stieltjes transfer]\label{lem:laplace}
Let $S(x)$ be a step function with locally finite variation, supported in
$[a,\infty)$ for some $a>0$. Suppose that, as $x\to\infty$,
\[
S(x)=P(\log x)+O(x^{-1}(1+\log x)^D)
\]
for a polynomial $P$. Interpret the Stieltjes integral below as an improper
integral if its total variation is not integrable. Then
\begin{equation}\label{eq:laplace}
\int_0^\infty e^{-\epsilon x}\dd S(x)
=\sum_{j=0}^{\deg P}\frac{\Gamma^{(j)}(1)}{j!}P^{(j)}(L)
 +O\!\left(\epsilon(1+L)^{D+1}\right),
\end{equation}
provided $D\ge\deg P$; an $o(1)$ remainder suffices without recording its rate.
\end{lemma}

\begin{proof}
Integration by parts gives $\epsilon\int_0^\infty e^{-\epsilon x}S(x)\dd x$.
The boundary terms vanish because $S=0$ near zero and grows at most
polynomially in $\log x$. Replacing $S$ by $P(\log x)$ on a bounded interval
costs $O(\epsilon)$, since powers of $\log x$ are integrable at zero. On
$x\ge A>1$, the error is bounded by
\[
C\epsilon\int_A^\infty e^{-\epsilon x}
 \frac{(1+\log x)^D}{x}\dd x
=O(\epsilon(1+L)^{D+1}).
\]
Split this integral at $x=1/\epsilon$ to obtain the last estimate.
For the main term put $t=\epsilon x$. The polynomial Taylor identity and the
Gamma integral give
\begin{align*}
\int_0^\infty e^{-t}P(L+\log t)\dd t
 &=\sum_{j=0}^{\deg P}\frac{P^{(j)}(L)}{j!}
     \int_0^\infty e^{-t}\log^j t\dd t\\
 &=\sum_{j=0}^{\deg P}\frac{\Gamma^{(j)}(1)}{j!}P^{(j)}(L).
\end{align*}
Only a finite sum is interchanged with the integral.
\end{proof}

\begin{theorem}[All-depth Abel endpoint polynomial]\label{thm:abel}
For every content $\mm$, define
\begin{equation}\label{eq:Q}
Q_{\mm,a}(L)=\Gamma(1+\partial_L)P_{\mm,a}(L)
:=\sum_{j=0}^{D(\mm)}\frac{\Gamma^{(j)}(1)}{j!}
                         \partial_L^jP_{\mm,a}(L).
\end{equation}
Then, for $d>0$,
\begin{equation}\label{eq:abel-asymptotic}
E_{\mm}(e^{-\epsilon};a)=Q_{\mm,a}(L)
 +O_{a,\mm}\bigl(\epsilon(1+L)^{D(\mm)+1}\bigr).
\end{equation}
The Abel endpoint constant is the ordinary limit after polynomial subtraction:
\begin{equation}\label{eq:Abelconstant}
A_{\mm}(a):=Q_{\mm,a}(0)
=\lim_{z\uparrow1}\bigl(E_{\mm}(z;a)
              -Q_{\mm,a}(-\log(1-z))+Q_{\mm,a}(0)\bigr).
\end{equation}
It belongs to the algebra of Theorem~\ref{thm:cutoff}, enlarged only by
$\gamma$ and ordinary zeta values $\zeta(2),\ldots,\zeta(D(\mm))$.
\end{theorem}

\begin{proof}
Let $S(x)$ sum the finite products over tuples whose largest lattice point is
at most $x$. Its jumps are exactly the coefficients of the Abel series in
\eqref{eq:E}. Between lattice points,
$S(x)=e_{\mm,N(x)}(a)$ with $N(x)+a=x+O(1)$. Thus
Theorem~\ref{thm:cutoff} gives the hypothesis of Lemma~\ref{lem:laplace}.
Applying that lemma proves~\eqref{eq:abel-asymptotic}. Replacing $L$ by
$-\log(1-z)$ changes only the vanishing error, proving the limit statement.
Finally,
\begin{equation}\label{eq:loggamma}
\log\Gamma(1+u)=-\gamma u+
 \sum_{k\ge2}\frac{(-1)^k\zeta(k)}{k}u^k
\end{equation}
shows that each $\Gamma^{(j)}(1)$ is a rational polynomial in the asserted
ordinary constants. The empty content is again handled by definition.
\end{proof}

The first Gamma moments are
\begin{align}
\Gamma'(1)&=-\gamma,\notag\\
\Gamma''(1)&=\gamma^2+\zeta(2),\notag\\
\Gamma'''(1)&=-\gamma^3-3\gamma\zeta(2)-2\zeta(3),\notag\\
\Gamma^{(4)}(1)&=\gamma^4+6\gamma^2\zeta(2)+3\zeta(2)^2
                +8\gamma\zeta(3)+6\zeta(4).
\label{eq:gamma-moments}
\end{align}
These formulas explain exactly which new ordinary constants an Abel endpoint
can introduce. No new positive-index generalized Stieltjes constant is
created by the transfer itself.

\subsection{A formal master series, and a necessary convergence warning}

Define, as a formal multivariate power series,
\begin{equation}\label{eq:K}
\KK(\uu)=\sum_{\mm\in\N_0^{R+1}}
 \frac{\Gamma^{(D(\mm))}(1)}
      {\prod_{r=0}^R m_r!(r+1)^{m_r}}\uu^{\mm}.
\end{equation}
Theorem~\ref{thm:abel} gives the compact master identity
\begin{equation}\label{eq:master}
\boxed{\displaystyle
\sum_{\mm}A_{\mm}(a)\uu^{\mm}=\CC_a(\uu)\KK(\uu).}
\end{equation}
Indeed, the coefficient of $\uu^{\mm}$ in $e^{V_{\uu}(\log t)}$ integrates
against $e^{-t}\dd t$ to the coefficient in~\eqref{eq:K}. For a fixed content
this is a finite, valid integration of logarithmic moments.

\begin{proposition}[The Gamma-moment master series need not converge]
\label{prop:formal}
When a $u_1$ variable is present, $\KK(0,u_1,0,\ldots)$ has radius of
convergence zero. Consequently~\eqref{eq:master} is not, in general, an
identity of analytic germs in all logarithmic-decoration variables.
\end{proposition}

\begin{proof}
The coefficient of $u_1^m$ is $\Gamma^{(2m)}(1)/(2^m m!)$. All integrands in
\[
\Gamma^{(2m)}(1)=\int_0^\infty e^{-t}\log^{2m}t\dd t
\ge e^{-1}\int_0^1\log^{2m}t\dd t=e^{-1}(2m)!
\]
are nonnegative. The $m$th roots of $(2m)!/(2^m m!)$ tend to infinity.
This proves the assertion. In particular, the tempting expression
$\int_0^\infty e^{-t}\exp(u_1\log^2t/2)\dd t$ diverges near $t=0$ for every
real $u_1>0$.
\end{proof}

This warning does not affect any fixed coefficient, nor the analytic
one-decoration generators below. It does prevent an invalid application of
holomorphic parameter differentiation to the unrestricted master integral.
The entire canonical product $\CC_a$ and the formal moment series $\KK$ have
fundamentally different convergence properties.

\subsection{A Gamma--polygamma differential recurrence}

On polynomials in $L$, put $D_L=\partial_L$, $\Gop=\Gamma(1+D_L)$, and
\begin{equation}\label{eq:Top}
\Top=L+\psi(1+D_L)
=L-\gamma+\sum_{k\ge1}(-1)^{k+1}\zeta(k+1)D_L^k.
\end{equation}
All operator series terminate on a fixed polynomial.

\begin{proposition}[Closed recurrence for endpoint polynomials]
\label{prop:operator}
With a negative content interpreted as zero,
\begin{equation}\label{eq:Q-recurrence}
\partial_LQ_{\mm,a}(L)=\sum_{r=0}^R\Top^r Q_{\mm-\ee_r,a}(L).
\end{equation}
The powers of $\Top$ mean composition of the displayed operator, not ordinary
commutative powers of its separate summands.
\end{proposition}

\begin{proof}
The relation $[D_L,L]=1$ gives
$\Gop L\Gop^{-1}=L+\psi(1+D_L)$. This follows first for a formal series
$G(D_L)$ from $G(D_L)L=LG(D_L)+G'(D_L)$; the constant term of $G$ is one here,
so its inverse exists on polynomials. Hence
$\Gop L^r\Gop^{-1}=\Top^r$. Differentiating
$\CC_a(\uu)e^{V_{\uu}(L)}$ gives
$\partial_LP_{\mm,a}=\sum_r L^rP_{\mm-\ee_r,a}$. Conjugate this equality by
$\Gop$, which commutes with $\partial_L$, to obtain~\eqref{eq:Q-recurrence}.
\end{proof}

\section{Gamma quotients, common shifts, and parameter differentiation}
\label{sec:unweighted}

When $u_0=u$ and the other auxiliary variables vanish, the canonical product
collapses to an inverse Gamma quotient:
\begin{equation}\label{eq:Cunweighted}
\CC_a(u,0,\ldots)=\frac{\Gamma(a)}{\Gamma(a+u)}.
\end{equation}
To see this, compare the logarithms near zero. The linear coefficient is
$-\psi(a)$, and the coefficient of $u^k$ for $k\ge2$ is
$(-1)^{k-1}\zeta(k,a)/k$, in agreement with~\eqref{eq:cumulants}.
Since both sides are entire functions of $u$, the local identity proves the
formula everywhere.

\begin{corollary}[Unweighted endpoint generator]\label{cor:unweighted}
For the depth-$d$ content $(d)$,
\begin{align}
\sum_{d\ge0}P_{(d),a}(L)u^d
 &=\frac{\Gamma(a)}{\Gamma(a+u)}e^{uL},\notag\\
\sum_{d\ge0}Q_{(d),a}(L)u^d
 &=\Acal_a(u)e^{uL},\qquad
\Acal_a(u):=\frac{\Gamma(a)\Gamma(1+u)}{\Gamma(a+u)}.
\label{eq:Aunweighted}
\end{align}
In particular, $A_{(d)}(a)=[u^d]\Acal_a(u)$. At $a=1$,
\begin{equation}\label{eq:aone}
Q_{(d),1}(L)=\frac{L^d}{d!},\qquad A_{(d)}(1)=0\quad(d\ge1).
\end{equation}
\end{corollary}

These endpoint identities have an independent full-function check. Let
$\mathcal E(z,u;a)=1+\sum_{d\ge1}u^dE_{(d)}(z;a)$.

\begin{proposition}[An exact incomplete-beta generator]\label{prop:beta}
For $a>0$, $0<z<1$, and $u$ near zero,
\begin{align}
\mathcal E(z,u;a)
 &=1+\frac{u z^a}{a}\,{}_2F_1(1,a+u;a+1;z)\label{eq:Ehyper}\\
 &=1+u(1-z)^{-u}B_z(a,u),\label{eq:Ebeta}
\end{align}
where $B_z(a,u)=\int_0^z t^{a-1}(1-t)^{u-1}\dd t$.
As $z\uparrow1$, locally uniformly in a small complex disc about $u=0$,
\begin{equation}\label{eq:beta-limit}
\mathcal E(z,u;a)=\Acal_a(u)(1-z)^{-u}+O(1-z).
\end{equation}
\end{proposition}

\begin{proof}
Condition on the largest index $n$. Its contribution to the generating
function is
\[
\frac{u z^{n+a}}{n+a}
 \prod_{j=0}^{n-1}\left(1+\frac{u}{a+j}\right)
=\frac{u z^a}{a}\frac{(a+u)_n}{(a+1)_n}z^n.
\]
Since $(1)_n/n!=1$, summing gives~\eqref{eq:Ehyper}. Euler's
hypergeometric transformation gives~\eqref{eq:Ebeta}; alternatively expand
its integral into a binomial series. The incomplete integral is entire in
$u$ for fixed $z<1$, so the identity is not limited to $\Rea u>0$.

Initially for $\Rea u>0$, write
$B_z(a,u)=B(a,u)-B_{1-z}(u,a)$. The complete beta term is
$uB(a,u)=\Acal_a(u)$. Taylor expansion of $(1-t)^{a-1}$ in the complementary
integral gives
\[
u(1-z)^{-u}B_{1-z}(u,a)=1+O(1-z).
\]
The coefficients in this last expansion extend holomorphically through
$u=0$, since their only possible denominators are $u+j$ with $j\ge1$ after
the constant term. This proves uniformity near zero and~\eqref{eq:beta-limit}.
\end{proof}

At the half shift,
\begin{equation}\label{eq:half-log}
\log\Acal_{1/2}(u)=2\log2\,u+
 \sum_{k\ge2}\frac{(-1)^k(2-2^k)\zeta(k)}k u^k.
\end{equation}
Consequently
\begin{align}
A_{(1)}(1/2)&=2\log2,\notag\\
A_{(2)}(1/2)&=2\log^22-\zeta(2),\notag\\
A_{(3)}(1/2)&=\frac43\log^32-2\zeta(2)\log2+2\zeta(3).
\label{eq:half-examples}
\end{align}
For a positive integer $M$,
\begin{equation}\label{eq:integer-A}
\Acal_M(u)=\prod_{j=1}^{M-1}\left(1+\frac uj\right)^{-1}.
\end{equation}
Its coefficient of $u^d$ is $(-1)^d$ times the degree-$d$ complete homogeneous
symmetric function of the reciprocals of the positive integers smaller than
$M$. Equivalently,
\begin{equation}\label{eq:integer-bell}
A_{(d)}(M)=\frac{(-1)^d}{d!}
 Y_d\bigl(H_{M-1},1!H_{M-1}^{(2)},\ldots,(d-1)!H_{M-1}^{(d)}\bigr),
\end{equation}
where the complete exponential Bell polynomials satisfy
\[
Y_d(x_1,\ldots,x_d)=d![t^d]\exp\left(\sum_{j=1}^d\frac{x_jt^j}{j!}\right).
\]
Thus this specialization gives finite generalized-harmonic identities at
all depths, not merely a collection of isolated low-order constants.

\subsection{Exact shift and derivative laws}

\begin{proposition}[Deletion of the first lattice point]\label{prop:shift}
As entire functions, with removable quotients interpreted by cancellation,
\begin{equation}\label{eq:Cshift}
\CC_{a+1}(\uu)=\frac{\CC_a(\uu)}{1+W_{\uu}(a)}.
\end{equation}
The same law holds for the formal Abel-constant generator in
\eqref{eq:master}. In coefficient form,
\begin{equation}\label{eq:Ashift}
A_{\mm}(a+1)+\sum_{r=0}^R f_r(a)A_{\mm-\ee_r}(a+1)=A_{\mm}(a).
\end{equation}
\end{proposition}

\begin{proof}
The Hurwitz recurrence gives
$\gamma_r(a+1)=\gamma_r(a)-f_r(a)$. Remove the $n=0$ factor from
\eqref{eq:Cproduct}; its exponential compensation cancels this Stieltjes
shift exactly. This proves~\eqref{eq:Cshift}. The formal factor $\KK$ is
independent of $a$, so multiplication by it proves the Abel law. Extraction
of a content gives~\eqref{eq:Ashift}.
\end{proof}

The elementary Hurwitz differential identity
$\partial_a\zeta(s,a)=-s\zeta(s+1,a)$ yields
\begin{equation}\label{eq:gamma-derivative}
\boxed{\displaystyle
\gamma_r'(a)=(-1)^{r+1}
 \bigl(\zeta^{(r)}(2,a)+r\zeta^{(r-1)}(2,a)\bigr),}
\end{equation}
where the second term is omitted for $r=0$. This follows by comparing the
coefficient of $t^r$ in $-(1+t)\zeta(2+t,a)$ and
\eqref{eq:stieltjes}. In particular,
\[
\gamma_0'(a)=-\zeta(2,a),\qquad
\gamma_1'(a)=\zeta'(2,a)+\zeta(2,a).
\]
Higher parameter derivatives use
\begin{equation}\label{eq:hurwitz-derivative}
\partial_a^q\zeta(s,a)=(-1)^q(s)_q\zeta(s+q,a).
\end{equation}
These identities differentiate every finite coefficient formula in this
article, producing further exact polygamma and zeta-derivative identities.
They do not confuse spectral differentiation with differentiation of the
Hurwitz parameter. Normal convergence of the canonical product on a small
complex neighborhood of a positive $a$ also justifies its parameter
 differentiation directly.

\section{One arbitrary logarithmic decoration at every depth}
\label{sec:one-decoration}

Let $\ell\ge1$. Write $P_{r;\ell,a}$, $Q_{r;\ell,a}$ and $A_{r;\ell}(a)$
for the content with $m_0=r$, $m_\ell=1$, and all other entries zero.
Here the total depth is $r+1$, not $r$. For instance,
\[
E_{r;\ell}(z;1)=\sum_{n_1>\cdots>n_{r+1}\ge1}
 \frac{z^{n_1}}{n_1\cdots n_{r+1}}
 \sum_{i=1}^{r+1}\log^\ell n_i.
\]
Define the analytic germ
\begin{align}
b_{a,\ell}(u)
&=\gamma_\ell(a)+\sum_{j\ge1}(-1)^{j+\ell}
                        \zeta^{(\ell)}(j+1,a)u^j\label{eq:bseries}\\
&=\gamma_\ell(a)-u\sum_{n\ge0}
 \frac{\log^\ell(n+a)}{(n+a)(n+a+u)}.
\label{eq:bresolvent}
\end{align}
The series in~\eqref{eq:bseries} converges for $|u|<a$; the normally
convergent resolvent sum gives its meromorphic continuation away from the
indicated lattice poles.

\begin{theorem}[Single-decoration identity and Stieltjes cancellation]
\label{thm:single-decoration}
For $u$ sufficiently close to zero,
\begin{equation}\label{eq:single-full}
\begin{split}
\sum_{r\ge0}Q_{r;\ell,a}(L)u^r
=\frac{\Gamma(a)}{\Gamma(a+u)}\biggl\{
 &b_{a,\ell}(u)\Gamma(1+u)e^{uL}\\
 &+\frac{1}{\ell+1}\partial_u^{\ell+1}
                \bigl(\Gamma(1+u)e^{uL}\bigr)\biggr\}.
\end{split}
\end{equation}
The derivative in braces acts only on its displayed Gamma--exponential
factor. In particular,
\begin{equation}\label{eq:single-constant}
\boxed{\displaystyle
\sum_{r\ge0}A_{r;\ell}(a)u^r
=\Acal_a(u)\left(b_{a,\ell}(u)
 +\frac{\Gamma^{(\ell+1)}(1+u)}{(\ell+1)\Gamma(1+u)}\right).}
\end{equation}
At $a=1$, for every $r\ge1$,
\begin{equation}\label{eq:single-cancellation}
\boxed{\displaystyle
A_{r;\ell}(1)=(-1)^{r+\ell}\zeta^{(\ell)}(r+1)
 +\frac1{\ell+1}[u^r]
             \frac{\Gamma^{(\ell+1)}(1+u)}{\Gamma(1+u)}.}
\end{equation}
Thus no positive-index Stieltjes constant occurs in this exact expression.
The depth-one exception is
\begin{equation}\label{eq:depth-one-decoration}
A_{0;\ell}(1)=\gamma_\ell+
                    \frac{\Gamma^{(\ell+1)}(1)}{\ell+1}.
\end{equation}
\end{theorem}

\begin{proof}
Use $u$ for the undecorated variable and $v$ for the single $\ell$th logarithm.
Differentiation of the logarithm of the canonical product at $v=0$ gives
\begin{align*}
\left.\partial_v\log\CC_a(u,v)\right|_{v=0}
 &=\gamma_\ell(a)+\sum_{n\ge0}f_\ell(n+a)
             \left(\frac1{1+u/(n+a)}-1\right)\\
 &=b_{a,\ell}(u).
\end{align*}
This differentiation is justified by normal convergence of the compensated
series. Expanding the resolvent proves~\eqref{eq:bseries}, including its sign.
The coefficient of $v$ in the cutoff generating function is therefore
\[
\frac{\Gamma(a)}{\Gamma(a+u)}e^{uL}
 \left(b_{a,\ell}(u)+\frac{L^{\ell+1}}{\ell+1}\right).
\]
Applying $\Gamma(1+\partial_L)$ and using
$L^{\ell+1}e^{uL}=\partial_u^{\ell+1}e^{uL}$ proves~\eqref{eq:single-full}.
The prefactor $\Gamma(a)/\Gamma(a+u)$ is not differentiated in this operation.
The equality can be read coefficientwise; equivalently the Gamma integral
is convergent for $u$ in a small disc, with only a fixed logarithmic power
inserted. Setting $L=0$ proves~\eqref{eq:single-constant}.

For $a=1$ the prefactor $\Acal_1(u)$ equals one exactly. The only
positive-index Stieltjes atom in $b_{1,\ell}$ is its constant term
$\gamma_\ell$. It cannot contribute to $[u^r]$ when $r\ge1$, giving
\eqref{eq:single-cancellation}. The remaining Gamma quotient is the complete
Bell polynomial
\[
\frac{\Gamma^{(\ell+1)}(1+u)}{\Gamma(1+u)}
=Y_{\ell+1}\bigl(\psi(1+u),\psi'(1+u),\ldots,\psi^{(\ell)}(1+u)\bigr).
\]
The Taylor coefficients of these polygamma functions involve only $\gamma$
and ordinary zeta values. This proves the asserted cancellation and the
depth-one exception.
\end{proof}

The cancellation is an algebraic consequence of the Gamma quotient, not an
assertion that any Stieltjes constant is independent of the surviving
constants. It is also stronger than a finite numerical observation: both the
depth $r+1$ and the logarithmic power $\ell$ are unrestricted.

\subsection{The single-logarithm family in ordinary zeta notation}

For $\ell=1$, equations~\eqref{eq:single-full}--\eqref{eq:single-constant}
become
\begin{align}
\sum_{r\ge0}Q_{r;1,a}(L)u^r
 &=\Acal_a(u)e^{uL}
  \left(b_{a,1}(u)+\frac{(L+\psi(1+u))^2+\psi'(1+u)}2\right),
  \label{eq:one-log-full}\\
\sum_{r\ge0}A_{r;1}(a)u^r
 &=\Acal_a(u)
  \left(b_{a,1}(u)+\frac{\psi(1+u)^2+\psi'(1+u)}2\right).
  \label{eq:one-log-constant}
\end{align}

\begin{corollary}[An explicit all-depth zeta-derivative identity]
\label{cor:one-log}
For every integer $r\ge1$,
\begin{equation}\label{eq:one-log-explicit}
\begin{split}
A_{r;1}(1)=(-1)^r\biggl(&-\zeta'(r+1)+\gamma\zeta(r+1)
 +\frac{r+1}{2}\zeta(r+2)\\
 &+\frac12\sum_{j=1}^{r-1}\zeta(j+1)\zeta(r-j+1)\biggr).
\end{split}
\end{equation}
The sum is empty for $r=1$.
\end{corollary}

\begin{proof}
In a disc about zero,
\[
\psi(1+u)=-\gamma+\sum_{j\ge1}(-1)^{j+1}\zeta(j+1)u^j,
\quad
\psi'(1+u)=\sum_{j\ge0}(-1)^j(j+1)\zeta(j+2)u^j.
\]
Multiply the first series by itself, add the second, divide by two, and
extract $[u^r]$. Add the coefficient
$(-1)^{r+1}\zeta'(r+1)$ from $b_{1,1}$. The result is
\eqref{eq:one-log-explicit}.
\end{proof}

The first three positive-depth-index constants are
\begin{align}
A_{1;1}(1)&=\zeta'(2)-\gamma\zeta(2)-\zeta(3),\label{eq:A11}\\
A_{2;1}(1)&=-\zeta'(3)+\gamma\zeta(3)+\frac{11}{4}\zeta(4),\label{eq:A21}\\
A_{3;1}(1)&=\zeta'(4)-\gamma\zeta(4)-2\zeta(5)-\zeta(2)\zeta(3).
\label{eq:A31}
\end{align}
Here $\zeta(2)^2=5\zeta(4)/2$ was used in~\eqref{eq:A21}.

For a completely explicit ordinary limit, put
\[
F(z)=\sum_{n>m\ge1}\frac{z^n(\log n+\log m)}{nm},
\qquad L_z=-\log(1-z).
\]
Then Theorem~\ref{thm:single-decoration} gives
\begin{equation}\label{eq:explicit-limit}
\begin{split}
\lim_{z\uparrow1}\biggl[
F(z)-\frac{L_z^3}{2}+\gamma L_z^2
-\left(\gamma_1+\frac{\gamma^2}{2}+\frac{3\zeta(2)}2\right)L_z
\biggr]
=\zeta'(2)-\gamma\zeta(2)-\zeta(3).
\end{split}
\end{equation}
Since $F(z)=-(\partial_{s_1}+\partial_{s_2})
\Li_{s_1,s_2}(z,1)|_{(1,1)}$, this is an exact endpoint identity for
polylogarithmic order derivatives. The constant on the right is not obtained
by assigning a value directly to the divergent series $F(1)$.

For a higher logarithmic power, the same theorem yields
\begin{equation}\label{eq:ell2}
A_{1;2}(1)=-\zeta''(2)+\gamma^2\zeta(2)+2\gamma\zeta(3)
                    +\frac92\zeta(4).
\end{equation}
The corresponding numerator is $\log^2 n+\log^2 m$, not
$\log n\log m$. To verify~\eqref{eq:ell2} directly from the generating
formula, differentiate
$\Gamma'''/\Gamma=\psi^3+3\psi\psi'+\psi''$ once at zero, divide by three,
and use~\eqref{eq:single-cancellation}. This produces
$\gamma^2\zeta(2)+\zeta(2)^2+2\gamma\zeta(3)+2\zeta(4)$,
which is the stated expression.

\section{Two logarithmic insertions and higher content}
\label{sec:two}

The notation $A_{(r,2)}$ in this section means $r$ undecorated factors and
\emph{two} factors carrying one logarithm each. It differs from
$A_{r;2}$, which denotes one factor carrying a squared logarithm.
Set $b_a=b_{a,1}$ and
\begin{equation}\label{eq:M}
M_a(u)=-\sum_{n\ge0}\frac{\log^2(n+a)}{(n+a+u)^2}
=\sum_{j\ge0}(-1)^{j+1}(j+1)\zeta''(j+2,a)u^j.
\end{equation}
The first sum is normally convergent off its lattice poles, and the Taylor
series is valid for $|u|<a$.

\begin{proposition}[Two-logarithm constant generator]\label{prop:two}
For $u$ near zero,
\begin{equation}\label{eq:two-generator}
\begin{split}
\sum_{r\ge0}A_{(r,2)}(a)u^r
=\Acal_a(u)\biggl[&\frac{b_a(u)^2+M_a(u)}2
 +\frac{b_a(u)\Gamma''(1+u)}{2\Gamma(1+u)}\\
 &+\frac{\Gamma^{(4)}(1+u)}{8\Gamma(1+u)}\biggr].
\end{split}
\end{equation}
\end{proposition}

\begin{proof}
The first and second derivatives of $\log\CC_a(u,v)$ at $v=0$ are
$b_a(u)$ and $M_a(u)$. Hence the coefficient of $v^2$ in $\CC_a(u,v)$ is
$\CC_a(u,0)(b_a^2+M_a)/2$. In the formal Gamma-moment factor, the
coefficients of $v^0,v^1,v^2$ are respectively
$\Gamma(1+u)$, $\Gamma''(1+u)/2$ and $\Gamma^{(4)}(1+u)/8$.
Multiply these three finite coefficient expressions and use
$\CC_a(u,0)\Gamma(1+u)=\Acal_a(u)$. This proves the identity without assuming
analyticity in the full $v$ variable.
\end{proof}

In particular,
\begin{equation}\label{eq:two-example}
\begin{split}
A_{(0,2)}(1)={}&\frac{\gamma_1^2-\zeta''(2)}2
 +\frac{\gamma_1(\gamma^2+\zeta(2))}2\\
 &+\frac{\gamma^4+6\gamma^2\zeta(2)+3\zeta(2)^2
                  +8\gamma\zeta(3)+6\zeta(4)}8.
\end{split}
\end{equation}
Unlike the single-decoration constants at depth greater than one, this
expression retains $\gamma_1$. The comparison with~\eqref{eq:ell2} is useful:
placing two logarithms on one index and placing them on two distinct indices
are analytically different operations, even though their total logarithmic
degree is the same.

The master formula gives every higher content by a finite recurrence. If
$\alpha\ne0$, $k=|\alpha|$, and $w=\sum_r r\alpha_r$, the coefficient of
$\uu^\alpha$ in $\log\CC_a$ is
\begin{equation}\label{eq:log-multi-coefficient}
\lambda_\alpha=
\begin{cases}
\gamma_r(a),&\alpha=\ee_r,\\[2pt]
\displaystyle
\frac{(-1)^{k-1+w}(k-1)!}{\prod_r\alpha_r!}\,
             \zeta^{(w)}(k,a),&k\ge2.
\end{cases}
\end{equation}
For cutoff polynomials replace the singleton coefficient by
$\gamma_r(a)+L^{r+1}/(r+1)$. If $F=\exp(\sum_{\alpha\ne0}
\lambda_\alpha\uu^\alpha)=\sum_\mm f_\mm\uu^\mm$, then
\begin{equation}\label{eq:euler-recurrence}
|\mm|f_\mm=\sum_{0<\alpha\le\mm}|\alpha|\lambda_\alpha
                                         f_{\mm-\alpha},
\qquad f_{\mathbf0}=1.
\end{equation}
This is obtained by applying $\sum_r u_r\partial_{u_r}$ to $F$.
It supplies an exact finite algorithm, with no symbolic infinite sums and no
integer-relation search.

\section{Diagonal spectral finite parts are a different operation}
\label{sec:diagonal}

For $\Rea s>1$ define
\[
Z_d(s;a)=\sum_{n_1>\cdots>n_d\ge0}\prod_{i=1}^d(n_i+a)^{-s},
\qquad Z_0(s;a)=1.
\]
Newton's identity gives the coefficientwise meromorphic identity
\begin{equation}\label{eq:diagonal-generator}
\sum_{d\ge0}Z_d(s;a)u^d
=\exp\!\left(\sum_{k\ge1}\frac{(-1)^{k-1}}k\zeta(ks,a)u^k\right).
\end{equation}
For any fixed $d$, only $k\le d$ is needed. Thus there is no need to define an
infinite exponential on a punctured parameter domain in order to continue
$Z_d$ near $s=1$.

Define the diagonal spectral constant
\begin{equation}\label{eq:Dd}
D_d(a)=\CT_{t=0}Z_d(1+t;a).
\end{equation}
Here $\CT$ means the constant coefficient of a one-variable Laurent series.
This is not, by definition, a cutoff or an Abel regularization.

\begin{theorem}[Exact diagonal anomaly]\label{thm:diagonal}
Let
\begin{equation}\label{eq:Rdiag}
\begin{split}
\RR_a(u,t)=\exp\biggl\{&u\left(\zeta(1+t,a)-\frac1t\right)\\
 &+\sum_{k\ge2}\frac{(-1)^{k-1}}k\zeta(k+kt,a)u^k\biggr\}.
\end{split}
\end{equation}
The expression is understood as a power series in $u$ with coefficients
holomorphic in $t$ near zero. Then
\begin{equation}\label{eq:diagonal-exact}
\boxed{\displaystyle
D_d(a)=\sum_{j=0}^d\frac1{j!}
               [u^{d-j}t^j]\RR_a(u,t).}
\end{equation}
The $j=0$ term is $C_{(d)}(a)$. The other terms give a finite explicit
correction involving Stieltjes constants only through index $d-1$ and
spectral derivatives of Hurwitz zeta at positive integers.
\end{theorem}

\begin{proof}
Equation~\eqref{eq:diagonal-generator} gives
$Z_d(1+t;a)=[u^d]e^{u/t}\RR_a(u,t)$. Expanding $e^{u/t}$ and taking the
constant Laurent coefficient gives~\eqref{eq:diagonal-exact}. At $t=0$, the
regular factor is exactly $\CC_a(u,0,\ldots)$, so the $j=0$ coefficient is the
cutoff constant. A Stieltjes coefficient of order $j$ is accompanied by at
least one $u$; hence the largest possible index in the extraction is $d-1$.
Everything is a finite coefficient calculation once $d$ is fixed.
\end{proof}

For $g_r=\gamma_r(a)$ the first cases are
\begin{align}
D_1(a)&=g_0,\notag\\
D_2(a)&=\frac{g_0^2-\zeta(2,a)}2-g_1,\label{eq:diag2}\\
D_3(a)&=\frac{g_0^3-3g_0\zeta(2,a)+2\zeta(3,a)}6
       -g_0g_1+\frac{g_2}{4}-\zeta'(2,a).
\label{eq:diag3}
\end{align}
For example, $Z_2(s;a)=(\zeta(s,a)^2-\zeta(2s,a))/2$. The product of the
pole $1/t$ with the linear coefficient $-g_1t$ contributes $-g_1$ to its
constant term. This contribution is absent in a cutoff constant extracted
before collapsing the spectral variables.

At $a=1$, the three unweighted schemes already differ in depth two:
\begin{equation}\label{eq:three-schemes}
\begin{array}{c|c}
\text{scheme and normalization}&\text{depth-two constant}\\ \hline
\text{sharp cutoff, polynomial in }\log N
  &(\gamma^2-\zeta(2))/2\\[2pt]
\text{Abel endpoint, polynomial in }-\log(1-z)&0\\[2pt]
\text{diagonal spectral Laurent constant}
  &(\gamma^2-\zeta(2))/2-\gamma_1
\end{array}
\end{equation}
These are different well-defined operations. In particular, the diagonal
constant must not be called the standard stuffle-regularized value without
an explicit change-of-regularization formula.

\subsection{Independent variables and general rays}

For absolutely convergent orders define the symmetrized shifted multiple
zeta function
\[
\mathscr Z(\mathbf s;a)=\sum_{\sigma\in S_d}
 \sum_{n_1>\cdots>n_d\ge0}
 \prod_{i=1}^d(n_i+a)^{-s_{\sigma(i)}}.
\]
The partition identity~\eqref{eq:partition} gives its meromorphic continuation
near $\mathbf s=\mathbf1$:
\begin{equation}\label{eq:symmetric-zeta}
\mathscr Z(\mathbf s;a)=\sum_{\pi\in\Pi_d}\mu(\pi)
                         \prod_{B\in\pi}\zeta\!\left(\sum_{i\in B}s_i,a\right).
\end{equation}
Only singleton blocks have a pole near this point. If the variables
$t_i=s_i-1$ remain independent, the multivariate constant coefficient of
$\prod_i(-\partial_{s_i})^{r_i}\mathscr Z$ equals the cutoff constant of the
labelled symmetric harmonic sum. This follows block by block: a singleton
has constant $\gamma_{r_i}(a)$, while a nonsingleton block has constant
$(-1)^{\sum_{i\in B}r_i}\zeta^{(\sum_{i\in B}r_i)}(|B|,a)$.
Since the blocks have disjoint variables, there is no hidden cross-block
Laurent contribution. Divide by $\prod_jm_j!$ to recover content normalization.

The situation changes upon restricting to a ray $t_i=c_it$, where $c_i\ne0$.
This restriction is made \emph{after} the independent spectral derivatives.
For a singleton with decoration $r_i$,
\begin{equation}\label{eq:singleton-ray}
(-1)^{r_i}\zeta^{(r_i)}(1+c_it,a)
=\frac{r_i!}{c_i^{r_i+1}t^{r_i+1}}
 +\sum_{m\ge0}\frac{(-1)^m\gamma_{r_i+m}(a)c_i^m}{m!}t^m.
\end{equation}
For a block $B$ of size at least two, put
$R_B=\sum_{i\in B}r_i$ and $c_B=\sum_{i\in B}c_i$.

\begin{proposition}[Finite formula on an arbitrary nonsingular ray]
\label{prop:ray}
For a partition $\pi$, let $S(\pi)$ be its singleton indices and put
$q_\pi=\sum_{i\in S(\pi)}(r_i+1)$. Define
\[
T_i(t)=\frac{r_i!}{c_i^{r_i+1}}
 +t^{r_i+1}\sum_{m\ge0}
       \frac{(-1)^m\gamma_{r_i+m}(a)c_i^m}{m!}t^m.
\]
The ray constant of the labelled differentiated symmetric sum is
\begin{equation}\label{eq:ray}
\sum_{\pi\in\Pi_d}\mu(\pi)[t^{q_\pi}]
 \left\{\prod_{i\in S(\pi)}T_i(t)
 \prod_{\substack{B\in\pi\\|B|\ge2}}
       (-1)^{R_B}\zeta^{(R_B)}(|B|+c_Bt,a)\right\}.
\end{equation}
Only Taylor coefficients up to $q_\pi$ are required in each summand.
\end{proposition}

\begin{proof}
In each term of~\eqref{eq:symmetric-zeta}, differentiation groups all
spectral derivatives belonging to a block into the derivative of that block's
single Hurwitz factor. Multiply each singleton expansion
\eqref{eq:singleton-ray} by $t^{r_i+1}$. The whole partition contribution has
then been multiplied by $t^{q_\pi}$. Its original Laurent constant is exactly
the coefficient displayed in~\eqref{eq:ray}.
\end{proof}

For example, the unweighted symmetric double has ray constant
\begin{equation}\label{eq:ray2}
g_0^2-\zeta(2,a)-\gamma_1(a)
                \left(\frac{c_1}{c_2}+\frac{c_2}{c_1}\right).
\end{equation}
Half of this at $c_1=c_2=1$ is~\eqref{eq:diag2}. Thus finite-part extraction
and restriction to a spectral diagonal do not commute. This elementary
counterexample is useful when moving an identity between a multivariate
Hurwitz-jet calculation and a one-variable spectral specialization.

\section{Exact antiderivatives and beta-weighted moment identities}
\label{sec:primitives}

This section gives convergent identities, with no endpoint finite-part
assignment. Two different integration variables are useful. Integration
against $\dd z/z$ raises the outer polylogarithmic index. Ordinary integration
against $\dd z$, or beta weights in $1-z$, has a different and particularly
simple all-depth structure.

\subsection{A kernel representation in the endpoint logarithm}

Put $L=-\log(1-z)$. Expanding~\eqref{eq:Ebeta} in $u$ and using
$v=-\log(1-t)$ gives, for every $d\ge1$,
\begin{equation}\label{eq:Lkernel}
\boxed{\displaystyle
E_{(d)}(z;a)=\frac1{(d-1)!}
 \int_0^L(L-v)^{d-1}(1-e^{-v})^{a-1}\dd v.}
\end{equation}
This identity also has a direct proof: the coefficient of $u^d$ in
$u(1-z)^{-u}B_z(a,u)$ is
\[
\frac1{(d-1)!}\int_0^z\frac{t^{a-1}}{1-t}
 \left(\log\frac{1-t}{1-z}\right)^{d-1}\dd t,
\]
and the stated change of variables yields~\eqref{eq:Lkernel}. Every integral
converges, since the kernel behaves as $v^{a-1}$ near zero.

Consequently $E_{(d)}(1-e^{-L};a)$ is the $d$-fold zero-based ordinary
antiderivative of $(1-e^{-L})^{a-1}$ in the $L$ variable. In particular,
$\partial_LE_{(d)}=E_{(d-1)}$ for $d\ge2$, while
$\partial_LE_{(1)}=(1-e^{-L})^{a-1}$. This is an exact full-function statement,
not merely an endpoint asymptotic.

\begin{theorem}[Beta-weighted all-depth moments]\label{thm:beta-moments}
For $a,b>0$ and integers $d\ge1$,
\begin{equation}\label{eq:beta-moment}
\boxed{\displaystyle
\int_0^1(1-z)^{b-1}E_{(d)}(z;a)\dd z
=\frac{B(a,b)}{b^d}.}
\end{equation}
For any integer $k\ge0$, one may differentiate $k$ times in $b$:
\begin{equation}\label{eq:log-beta-moment}
\int_0^1(1-z)^{b-1}\log^k(1-z)E_{(d)}(z;a)\dd z
=\frac{B(a,b)}{b^d}Y_k(\Delta_1,\ldots,\Delta_k),
\end{equation}
where $Y_0=1$ and
\begin{equation}\label{eq:Delta}
\Delta_j=\psi^{(j-1)}(b)-\psi^{(j-1)}(a+b)
                  +\frac{(-1)^j d(j-1)!}{b^j},\qquad j\ge1.
\end{equation}
\end{theorem}

\begin{proof}
With $L=-\log(1-z)$ the left side of~\eqref{eq:beta-moment} is
$\int_0^\infty e^{-bL}E_{(d)}(1-e^{-L};a)\dd L$.
The integrand in~\eqref{eq:Lkernel} is nonnegative, so Tonelli's theorem
allows the integration order to be changed. The inner Gamma integral is
\[
\frac1{(d-1)!}\int_v^\infty e^{-bL}(L-v)^{d-1}\dd L
=\frac{e^{-bv}}{b^d}.
\]
The remaining integral, after $t=e^{-v}$, is
$\int_0^1t^{b-1}(1-t)^{a-1}\dd t=B(a,b)$. This proves the moment formula.
The integral and every fixed $b$ derivative converge uniformly for $b$ in a
compact positive interval: the added factor is a power of $L$, while the
Laplace damping remains exponential. Hence differentiation under the
integral is valid. The logarithm of the right-hand side has $j$th derivative
$\Delta_j$, so the complete Bell-polynomial derivative rule gives
\eqref{eq:log-beta-moment}.
\end{proof}

Several particularly transparent identities follow:
\begin{align}
\int_0^1 E_{(d)}(z;a)\dd z&=\frac1a,\qquad d\ge1,\label{eq:moment-simple}\\
\frac1{(q-1)!}\int_0^1(1-z)^{q-1}E_{(d)}(z;a)\dd z
 &=\frac1{(a)_q q^d},\qquad q\in\N,\label{eq:qordinary}\\
\int_0^1 z^{h-1}E_{(d)}(z;a)\dd z
 &=\sum_{j=0}^{h-1}(-1)^j\binom{h-1}{j}
                   \frac{B(a,j+1)}{(j+1)^d},\qquad h\in\N.
\label{eq:positive-moments}
\end{align}
The last identity follows by expanding $z^{h-1}=(1-(1-z))^{h-1}$.
It is a finite exact evaluation, not an asymptotic expansion. For example,
\[
\int_0^1zE_{(d)}(z;a)\dd z
=\frac1a-\frac1{a(a+1)2^d}.
\]
A reflection specialization is especially simple. For $0<a<1$,
\begin{equation}\label{eq:reflection-moment}
\int_0^1(1-z)^{-a}E_{(d)}(z;a)\dd z
=\frac{\pi}{\sin(\pi a)(1-a)^d}.
\end{equation}
For example, put $\ell_2=\log2$. At the half lattice,
\begin{align}
\int_0^1\frac{E_{(d)}(z;1/2)}{\sqrt{1-z}}\dd z
 &=2^d\pi,\notag\\
\int_0^1\frac{\log(1-z)E_{(d)}(z;1/2)}{\sqrt{1-z}}\dd z
 &=-2^{d+1}\pi(d+\ell_2),\label{eq:half-moments}\\
\int_0^1\frac{\log^2(1-z)E_{(d)}(z;1/2)}{\sqrt{1-z}}\dd z
 &=2^d\pi\left(4(d+\ell_2)^2+4d+\frac{\pi^2}{3}\right).\notag
\end{align}
These follow from $B(1/2,1/2)=\pi$,
$\psi(1/2)-\psi(1)=-2\log2$, and
$\psi'(1/2)-\psi'(1)=\pi^2/3$. The family is uniform in the depth; the
logarithms here decorate the integration weight, and should not be confused
with spectral logarithms at the lattice indices.

For positive integers $a=M$, $b=N$, formula~\eqref{eq:Delta} becomes
\begin{equation}\label{eq:Delta-harmonic}
\Delta_j=(-1)^j(j-1)!
 \left(H_M^{(j)}(N)+\frac d{N^j}\right).
\end{equation}
Thus every logarithmic moment in~\eqref{eq:log-beta-moment} at such parameters
is an explicit rational expression in finite generalized harmonic numbers,
with a rational beta prefactor. The simple moment laws are special-function
consequences of the kernel representation; no priority is claimed for their
unshifted log-power specializations.

\begin{corollary}[Gamma-weighted generalized harmonic sums]
\label{cor:gamma-harmonic}
For $a,b>0$ and $d\ge1$,
\begin{equation}\label{eq:gamma-harmonic}
\sum_{n\ge0}e_{(d-1),n}(a)
          \frac{\Gamma(n+a)}{\Gamma(n+a+b+1)}
=\frac{\Gamma(a)}{b^d\Gamma(a+b)}.
\end{equation}
In particular,
\begin{align}
\sum_{n\ge0}\frac{\Gamma(n+a)H_n^{(1)}(a)}{\Gamma(n+a+b+1)}
 &=\frac{\Gamma(a)}{b^2\Gamma(a+b)},\notag\\
\sum_{n\ge0}\frac{\Gamma(n+a)
      \bigl((H_n^{(1)}(a))^2-H_n^{(2)}(a)\bigr)}{\Gamma(n+a+b+1)}
 &=\frac{2\Gamma(a)}{b^3\Gamma(a+b)}.
\label{eq:gamma-harmonic-examples}
\end{align}
\end{corollary}
\begin{proof}
Condition on the largest index in the beta moment. Tonelli's theorem gives
\[
\int_0^1(1-z)^{b-1}E_{(d)}(z;a)\dd z
=\sum_{n\ge0}\frac{e_{(d-1),n}(a)}{n+a}B(n+a+1,b).
\]
Write each beta factor as a Gamma quotient, use
Theorem~\ref{thm:beta-moments}, and cancel $\Gamma(b)$. The two displayed
examples use $e_{(1),n}=H_n^{(1)}$ and
$2e_{(2),n}=(H_n^{(1)})^2-H_n^{(2)}$.
\end{proof}

For an explicit half-lattice specialization, let
$\mathrm{Cat}_n=\binom{2n}{n}/(n+1)$ be the Catalan numbers. Since
\[
\frac{\Gamma(n+1/2)}{\Gamma(n+2)}
=\sqrt\pi\,\frac{\mathrm{Cat}_n}{4^n},\qquad
H_n^{(j)}(1/2)=2^jH_{2n}^{(j)}-H_n^{(j)},
\]
we obtain the convergent harmonic identities
\begin{align}
\sum_{n\ge0}\frac{\mathrm{Cat}_n}{4^n}(2H_{2n}-H_n)&=4,\label{eq:catalan1}\\
\sum_{n\ge0}\frac{\mathrm{Cat}_n}{4^n}
 \bigl((2H_{2n}-H_n)^2-4H_{2n}^{(2)}+H_n^{(2)}\bigr)&=16.
\label{eq:catalan2}
\end{align}
More generally, the corresponding sum of $e_{(d-1),n}(1/2)$ is $2^d$.
These identities are beta/Gauss consequences of the unweighted generator,
included to make the finite generalized-harmonic content explicit; no claim
that the individual low-depth Catalan specializations are previously unknown
is intended.

\subsection{Ordinary zero-based primitives and integer-shift elementary forms}

Define the $q$-fold ordinary primitive by
\[
I_{q,d}(z;a)=\frac1{(q-1)!}\int_0^z(z-t)^{q-1}E_{(d)}(t;a)\dd t,
\qquad q\ge1.
\]
Termwise integration of~\eqref{eq:Ehyper} gives
\begin{equation}\label{eq:ordinary-primitive-generator}
\sum_{d\ge1}I_{q,d}(z;a)u^d
=\frac{u\Gamma(a)z^{a+q}}{\Gamma(a+q+1)}
          {}_2F_1(1,a+u;a+q+1;z).
\end{equation}
At $z=1$, Gauss summation reduces the right-hand side to
$u/((a)_q(q-u))$, agreeing with~\eqref{eq:qordinary} coefficient by
coefficient. Either the beta-moment proof or the convergent Gauss formula
justifies this endpoint evaluation for $u$ near zero.

At a positive integer shift the kernel in~\eqref{eq:Lkernel} is a finite
exponential polynomial. Therefore the full nested harmonic function is
elementary:
\begin{equation}\label{eq:integer-elementary}
\begin{split}
E_{(d)}(z;M)=\frac{L^d}{d!}
+\sum_{j=1}^{M-1}\frac{(-1)^j\binom{M-1}{j}}{(-j)^d}
 \left((1-z)^j-\sum_{k=0}^{d-1}\frac{(-jL)^k}{k!}\right),
\quad L=-\log(1-z).
\end{split}
\end{equation}
To prove this, expand
$(1-e^{-v})^{M-1}=\sum_{j=0}^{M-1}(-1)^j\binom{M-1}{j}e^{-jv}$ in
\eqref{eq:Lkernel} and integrate each exponential $d$ times with zero initial
conditions. The $j=0$ term gives $L^d/d!$; the Taylor subtraction in the
other terms enforces those initial conditions exactly.

For $M=2$ the first two cases are
\[
E_{(1)}(z;2)=L-z,\qquad E_{(2)}(z;2)=\frac{L^2}{2}-L+z.
\]
Dropping the exponentially small $(1-z)^j$ terms in
\eqref{eq:integer-elementary} gives precisely the endpoint polynomials
from~\eqref{eq:integer-A}. Equating their constant coefficients also recovers
the finite symmetric-function identity
\[
h_d(1,1/2,\ldots,1/(M-1))
=\sum_{j=1}^{M-1}\frac{(-1)^{j-1}\binom{M-1}{j}}{j^d},
\qquad M\ge2,
\]
where $h_d$ is the complete homogeneous symmetric function. This finite
identity is another classical specialization, not an independence statement.

\subsection{Logarithmic primitives raise the outer polylogarithmic index}

Let $D_z=z\partial_z$. For a nonempty content and $p\ge1$, put
\begin{equation}\label{eq:Jdef}
J_{p,\mm}(z;a)=\frac1{(p-1)!}
 \int_0^z\left(\log\frac zt\right)^{p-1}E_{\mm}(t;a)\frac{\dd t}{t}.
\end{equation}
Then $D_zJ_{p,\mm}=J_{p-1,\mm}$, where $J_{0,\mm}=E_{\mm}$, and all these
primitives vanish at $z=0$. Termwise integration gives
\begin{equation}\label{eq:Jseries}
J_{p,\mm}(z;a)=\sum_{n_1>\cdots>n_d\ge0}
 \frac{z^{n_1+a}}{(n_1+a)^p}
 \sum_{\cont(\mathbf r)=\mm}\prod_i f_{r_i}(n_i+a).
\end{equation}
For $z=1$ this is absolutely convergent. To see this, sum the absolute inner
factors first. If the outer decoration is $r$, the inner coefficient is
$O((1+\log n)^{D-r-1})$ by elementary bounds on harmonic sums. The outer
factor is therefore $O((1+\log n)^{D-1}/n^{p+1})$, a summable sequence.
This also justifies the endpoint integral in~\eqref{eq:Jdef}.

For unweighted content $(d)$, the full generating function is
\begin{equation}\label{eq:Jhyper}
\sum_{d\ge1}u^dJ_{p,(d)}(z;a)
=\frac{u z^a}{a^{p+1}}
 {}_{p+2}F_{p+1}\!\left(
 \begin{matrix}1,a+u,\underbrace{a,\ldots,a}_{p}\\
               \underbrace{a+1,\ldots,a+1}_{p+1}\end{matrix};z\right).
\end{equation}
Indeed, condition on the largest lattice point as in
Proposition~\ref{prop:beta}, and multiply its coefficient by $(n+a)^{-p}$.
The ratio $(a)_n/(a+1)_n=a/(a+n)$ supplies each extra denominator.
At $z=1$ the hypergeometric excess is $p-u$, so the expression is convergent
for $\Rea u<p$ in its usual parameter domain; a neighborhood of zero is enough
for every coefficient identity used here.

\subsection{All positive integer shifts: a finite Gamma deletion theorem}

For $M,p,d\in\N$, abbreviate
\[
J_{p,d}(M)=J_{p,(d)}(1;M)
=\sum_{n_1>\cdots>n_d\ge M}
 \frac1{n_1^{p+1}n_2\cdots n_d},\qquad
P_{M-1}(u)=\prod_{j=1}^{M-1}\left(1+\frac uj\right).
\]

\begin{theorem}[Integer-shift convergent antiderivative hierarchy]
\label{thm:integer-primitives}
As an identity of analytic germs at $(u,v)=(0,0)$,
\begin{equation}\label{eq:integer-J}
\boxed{\displaystyle
\sum_{p,d\ge1}J_{p,d}(M)v^p u^d
=\frac{1}{P_{M-1}(u)}
 \left(1-\frac{\Gamma(1-u)\Gamma(1-v)}{\Gamma(1-u-v)}
       -\sum_{n=1}^{M-1}\frac{(u)_n}{n!}\frac{v}{n-v}\right).}
\end{equation}
Every coefficient belongs to the $\Q$-algebra generated by ordinary zeta
values. In particular,
\begin{equation}\label{eq:first-J}
\sum_{d\ge1}J_{1,d}(M)u^d
=\frac{-\psi(1-u)-\gamma
             -\sum_{n=1}^{M-1}(u)_n/(n!n)}{P_{M-1}(u)}.
\end{equation}
\end{theorem}

\begin{proof}
The generating coefficient conditioned on the largest integer $n\ge M$ is
\[
\sum_{d\ge1}\frac{u^d}{n^{p+1}}
 e_{d-1}(1/M,\ldots,1/(n-1))
=\frac1{P_{M-1}(u)}\frac{(u)_n}{n!\,n^p}.
\]
Thus summing $p$ gives the absolutely convergent series, locally near zero,
\begin{equation}\label{eq:Jresolvent}
\frac1{P_{M-1}(u)}\sum_{n\ge M}\frac{(u)_n}{n!}\frac{v}{n-v}.
\end{equation}
For example $|u|,|v|<1/4$ suffices: the summands decay as
$O(n^{-2+\Rea u})$ uniformly on a smaller compact set.

For $M=1$ use the binomial identity
$\sum_{n\ge1}(u)_nt^n/n!=(1-t)^{-u}-1$. Initially take $\Rea v<0$,
$\Rea u<1$. Integration gives
\begin{align*}
\sum_{n\ge1}\frac{(u)_n}{n!}\frac{v}{n-v}
 &=v\int_0^1t^{-v-1}\bigl((1-t)^{-u}-1\bigr)\dd t\\
 &=vB(-v,1-u)+1\\
 &=1-\frac{\Gamma(1-u)\Gamma(1-v)}{\Gamma(1-u-v)}.
\end{align*}
Both sides continue holomorphically near $(0,0)$, so this proves the classical
height-one Gamma identity there. Remove the finitely many terms $n<M$ in
\eqref{eq:Jresolvent} to obtain~\eqref{eq:integer-J}. The Gamma logarithm
\eqref{eq:loggamma} shows that its coefficients involve only ordinary zeta
values and rational numbers; the finite sum and $P_{M-1}^{-1}$ add only
rational coefficients. Finally extract the coefficient of $v$ to obtain
\eqref{eq:first-J}.
\end{proof}

For $M=1$, equation~\eqref{eq:first-J} is the classical height-one identity
$J_{1,d}(1)=\zeta(d+1)$. At $M=2$ it gives
\begin{align}
J_{1,1}(2)&=\zeta(2)-1,\notag\\
J_{1,2}(2)&=\zeta(3)-\zeta(2)+1,\notag\\
J_{1,3}(2)&=\zeta(4)-\zeta(3)+\zeta(2)-1.
\label{eq:Jexamples}
\end{align}
The theorem covers all $p$ and $d$, not only this first-primitive line.
It does not assert that the general rational-shift hypergeometric expression
\eqref{eq:Jhyper}, or its logarithmically decorated analogues, reduces to the
same single-zeta algebra.

\section{Rational shifts and cyclotomic transport}
\label{sec:cyclotomic}

Let $a=h/q$ with $1\le h\le q$, and write $\omega=e^{2\pi i/q}$ and
$w=z^{1/q}$ for $0<z<1$. Define the common-shift ordered polylogarithmic sum
\[
\mathcal F_{\mathbf s}(z;a)
=\sum_{n_1>\cdots>n_d\ge0}z^{n_1+a}
                       \prod_i(n_i+a)^{-s_i}.
\]

\begin{proposition}[Root-of-unity filter with its normalization]
\label{prop:filter}
For $|z|<1$ on a fixed root branch, and in particular on the positive real
interval,
\begin{equation}\label{eq:filter}
\begin{split}
\mathcal F_{\mathbf s}(z;h/q)
=q^{s_1+\cdots+s_d-d}
 \sum_{j_1,\ldots,j_d=0}^{q-1}\omega^{-h(j_1+\cdots+j_d)}
 \Li_{\mathbf s}(w\omega^{j_1},\omega^{j_2},\ldots,\omega^{j_d}).
\end{split}
\end{equation}
The identity may be spectrally differentiated to any fixed order. In doing
so, the factor $q^{s_1+\cdots+s_d-d}$ must also be differentiated.
\end{proposition}

\begin{proof}
Put $k_i=qn_i+h$. All $k_i$ are positive, strictly ordered and congruent to
$h$ modulo $q$. The denominators contribute $q^{s_1+\cdots+s_d}$, and the
outer exponential is $w^{k_1}$. Insert the exact filter
\[
\mathbf1_{k\equiv h\ (q)}=\frac1q\sum_{j=0}^{q-1}\omega^{j(k-h)}
\]
for each index and interchange the finite sums. This gives~\eqref{eq:filter}.
Exponential damping makes the underlying sums and all fixed spectral
derivatives locally normally convergent. Differentiation of the prefactor
is necessary because $\log(n+a)=\log k-\log q$.
\end{proof}

At the unweighted orders $s_i=1$, the prefactor is one ($q^0$). For the logarithmic primitive with $s_1=p+1$ and the other
orders one, it is $q^p$. The latter formula remains absolutely convergent at
$z=1$ for $p\ge1$, even after fixed spectral differentiations. Thus the
convergent antiderivative results provide exact finite combinations of
cyclotomic multiple-polylogarithm values, without any finite-part ambiguity.
The beta moments similarly translate to identities involving root-of-unity
polylogarithms and the Gamma ratio $B(h/q,b)$.

\subsection{The logarithmic scale changes under root extraction}

A separate adjustment is needed at a divergent Abel endpoint. With
$z=e^{-\epsilon}$ and $w=e^{-\epsilon/q}$, the logarithmic endpoint variables
are
\[
L=\log(1/\epsilon),\qquad T=\log(q/\epsilon)=L+\log q.
\]
After spectral differentiation, let the \emph{complete filtered sum} on the
right of~\eqref{eq:filter} represent $E_{\mm}(z;h/q)$. Its polynomial in $T$
is therefore
\begin{equation}\label{eq:scale}
Q_{\mm,h/q}(T-\log q).
\end{equation}
The constant in the $T$ normalization is $Q_{\mm,h/q}(-\log q)$, whereas the
constant in the original $L$ normalization is $Q_{\mm,h/q}(0)$. Their exact
difference is
\begin{equation}\label{eq:scale-difference}
Q_{\mm,h/q}(-\log q)-Q_{\mm,h/q}(0)
=\sum_{j=1}^{D(\mm)}\frac{(-\log q)^j}{j!}Q_{\mm,h/q}^{(j)}(0).
\end{equation}
This statement concerns the full finite trace and does not require a claim
about independent asymptotic expansions of its individual colored summands.

For the smallest example, $a=1/2$ gives
\[
E_{(1)}(z;1/2)=\Li_1(\sqrt z)-\Li_1(-\sqrt z)
=\log\frac{1+\sqrt z}{1-\sqrt z}.
\]
Its constant relative to $-\log(1-z)$ is $2\log2$. Its constant relative to
$-\log(1-\sqrt z)$ is only $\log2$. The difference is not an error in either
polylogarithm identity; it is a change in the specified endpoint scale.
A cyclotomic coefficient compiler must preserve this distinction.

\section{Reproducible coefficient computations and validation}
\label{sec:verification}

The package implements the finite recurrence~\eqref{eq:euler-recurrence}, not
an unbounded symbolic expansion. A requested content bounds every auxiliary
variable separately. The cutoff polynomial is computed first; the finite
Gamma differential operator then gives the Abel polynomial. The same
constant is independently recomputed by convolution with~\eqref{eq:K}.

The symbols in the implementation have the following meanings:
\begin{center}
\begin{tabular}{ll}
\toprule
Symbol & Meaning\\
\midrule
\texttt{g0}, \texttt{g1}, \ldots & $\gamma_0(a),\gamma_1(a),\ldots$\\
\texttt{z2\_1}, \texttt{z3\_2}, \ldots
 & $\zeta'(2,a),\zeta''(3,a),\ldots$\\
\texttt{G} & Euler's constant $\gamma$\\
\texttt{z2}, \texttt{z3}, \ldots & Ordinary zeta values\\
\texttt{L} & The explicitly specified endpoint logarithm\\
\bottomrule
\end{tabular}
\end{center}
Treating these atoms as indeterminates is a device for exact polynomial
verification. It does not assume that the corresponding numbers are
algebraically or rationally independent.

\subsection{Finite exact arithmetic}

The exact suite compares direct finite products on rational alphabets with
power-sum recurrences;
checks the labelled set-partition formula; compares the two implementations
of Gamma transfer; verifies the unweighted cancellation, single- and
double-decoration generators, higher single decorations, and the
Gamma-conjugated differential recurrence; compares Newton and partition
formulas for diagonal constants; verifies half and integer shifts; checks
finite deletion identities for primitives and integer-shift elementary
endpoint polynomials; and checks root-of-unity filters by exact remainders
modulo cyclotomic polynomials.

The executed exact suite passed 2,600 assertions. Their grouping is recorded
in \src{results/exact_checks.json} and summarized in the accompanying README. These are finite regression tests of
the formulas and implementation. They do not replace the analytic proofs of
normal convergence, asymptotic transfer, or interchange of infinite sums and
integrals. No proof assistant was run.

\subsection{Independent numerical checks}

The executed numerical suite passed 158 comparisons at 60 decimal digits.
The largest observed absolute residual was approximately
$7.83\times10^{-55}$; full values and tolerances are retained in
\src{results/numerical_checks.json}. It compares finite decorated
harmonic products with Stieltjes differences and Hurwitz-zeta derivatives;
direct geometric-convergent sums with the hypergeometric and incomplete-beta
generators; numerical integrals with the convergent primitive and beta-moment
formulas; numerical parameter derivatives of Stieltjes constants with
\eqref{eq:gamma-derivative}; and Cauchy averages of Laurent functions with the
diagonal constants. It also checks the integer-shift elementary functions
against direct harmonic coefficients.

These are ordinary floating-point computations, not interval enclosures.
Reported residuals are observed residuals, not rigorous error bounds. In
particular, a very small residual would not prove an identity that lacked
its written proof. The exact suite and the numerical suite deliberately use
different formulations where practical.

\subsection{Direct endpoint diagnostics}

A separate script sums the decorated Abel series directly and subtracts the
\emph{entire} predicted polynomial. It uses $N=\lceil35/\epsilon\rceil$ terms,
with $\epsilon=10^{-2},10^{-3},10^{-4},10^{-5}$, for two shifts and four
contents. At the smallest $\epsilon$ each sum uses 3,500,000 terms.
The following observed differences at $a=1$ illustrate why naive endpoint
extrapolation can be misleading:
\begin{center}
\begin{tabular}{rcc}
\toprule
$\epsilon$ & $E_{(1,1)}-Q_{(1,1),1}$ & $E_{(1,2)}-Q_{(1,2),1}$\\
\midrule
$10^{-2}$ & $1.6344939\cdot10^{-1}$ & $1.7682030$\\
$10^{-3}$ & $3.0587602\cdot10^{-2}$ & $7.0392944\cdot10^{-1}$\\
$10^{-4}$ & $5.0120208\cdot10^{-3}$ & $2.0368075\cdot10^{-1}$\\
$10^{-5}$ & $7.4952623\cdot10^{-4}$ & $4.7670923\cdot10^{-2}$\\
\bottomrule
\end{tabular}
\end{center}
Here $Q$ is evaluated at $L=\log(1/\epsilon)$. The displayed differences
contain the analytic remainder, finite truncation, and floating-point
roundoff; they are not certificates for any of those components separately.
Even millions of terms do not by themselves reveal many digits of the
constant in a higher logarithmic content. The coefficient formulas, rather
than extrapolated decimal values, determine the exact answers.

\section{Targeted corrections and integration guidance}
\label{sec:audit}

\subsection{A persistent independence wording issue}

In the reviewed snapshot, the canonical integration chapter, at the paragraph
labelled \src{integral:neg:psim2}, describes a dependent PSLQ input basket and
then says: ``Basis atoms must be $\Q$-independent.'' The numerical diagnosis is
useful, but the unrestricted requirement is misleading. Proving the needed
independence is often a difficult open problem, and it is neither an input
requirement for PSLQ nor an acceptable unstated assumption.

A precise replacement is:
\begin{quote}
Remove known exact rational dependencies from the candidate list, or work
with representatives modulo the explicitly proved relation system. A
relation finder may otherwise return an already-known relation that does not
involve the target. No unproved rational independence of the remaining
constants is assumed, and failure to find a target relation proves neither
independence nor nonreducibility.
\end{quote}

This is a correction to persistent wording, not a claim to have newly
discovered the general PSLQ caveat. The manuscript's broader distinction
between numerical evidence and proof is retained. A replacement paragraph
is supplied as an editorial fragment, not as an automatically applied patch.

\subsection{Regularization conventions that should be stated explicitly}

For any new endpoint or spectral identity, specify the limiting parameter,
the logarithmic variable, and the subtraction polynomial. The explicit
counterexample~\eqref{eq:three-schemes} shows why the bare phrase ``finite
part'' is insufficient when changing between cutoffs, Abel damping, and a
spectral diagonal. The scale law~\eqref{eq:scale-difference} supplies a second
necessary convention when transporting to roots of unity.

These are proposed convention improvements and safeguards for integrating
this continuation. We do \emph{not} assert that a particular unreviewed
incoming manuscript committed either mistake. Similarly,
Proposition~\ref{prop:formal} is a warning about this article's own generating
series, not an allegation that the repository already treated it as
holomorphic.

\subsection{Recommended placement and status}

A suitable self-contained report location is
\begin{center}
\src{Analysis/Polylogarithms/docs/reports/endpoint-regularization/}.
\end{center}
The harmonic-depth discussion can point to the decorated coefficient
calculus, while the integration and differentiation chapters can point to
the Gamma-transfer, Stieltjes-derivative and primitive results. The report
can also be cross-linked to the existing alternating-harmonic continuation,
but their normalizations and research questions should not be merged merely
because both use elementary symmetric harmonic functions.

The delivered files preserve a standalone article, proof status, runnable
code and recorded outputs. The canonical manuscript should be edited only
after independent mathematical review and comparison with the unexamined
incoming payloads. This article proves the families stated here. It does not
prove the repository's $S_6$ or $S_8$ candidates, settle period independence,
or establish minimum depth for any individual numerical value.

\section{Further research questions}
\label{sec:questions}

\subsection{Ordered rather than fully symmetrized logarithmic values}

The first obstruction beyond our theorem is already visible in depth two:
separate the endpoint constants with numerators $\log n$ and $\log m$ rather
than their sum. The present results determine their symmetric combination
but not their difference. Determine a functional or integral representation
for that antisymmetric component, then identify an exact basis of special
functions sufficient for its endpoint constant. This is a reduction question,
not a conjecture of arithmetic independence.

\subsection{Convergent spectral decorations of the primitive hierarchy}

Theorem~\ref{thm:integer-primitives} settles the unweighted integer-shift
family at all outer indices and depths. With a logarithmic insertion, the
series remains absolutely convergent after one $\dd z/z$ integration, but
the finite Gamma quotient no longer immediately evaluates the whole
hierarchy. Determine which symmetrized insertions arise from derivatives of
a finite-parameter Gamma or hypergeometric identity. The first target is the
one-logarithm primitive at arbitrary depth, with the outer and inner
spectral derivatives kept separate until the final symmetrization.

\subsection{Mixed shifts and rational lattice relations}

The common lattice is essential to the partition reduction. If each labelled
slot has its own shift $a_i$, collision blocks involve products of distinct
linear denominators rather than a single Hurwitz-zeta derivative. Partial
fractions should give a finite extension for distinct shifts, with confluent
limits producing polygamma derivatives. Establish a uniform formula that
remains regular as shifts collide. For rational common shifts, use
\eqref{eq:filter} and~\eqref{eq:scale-difference} to build explicit cyclotomic
endpoint relations at small conductors without silently changing the finite
part.

\subsection{Which combinations eliminate positive-index Stieltjes atoms?}

Theorem~\ref{thm:single-decoration} explains a complete cancellation for one
logarithmic decoration at $a=1$. Two distinct logarithmic insertions retain
$\gamma_1$ in~\eqref{eq:two-example}. For a fixed total logarithmic degree,
classify rational combinations of content constants whose coefficient
polynomials eliminate all positive-index Stieltjes symbols. This can first
be solved in the explicitly stated formal atom ring. Any subsequent claim
about the independence of numerical periods would require additional proof.

\subsection{Ray dependence and compatible renormalizations}

Formula~\eqref{eq:ray} reduces direction dependence to finitely many
coefficients at each depth. Classify rays, or finite linear combinations of
ray constants, that agree with independent-variable cutoff extraction
through a prescribed depth. Depth two gives the explicit ratio
$c_1/c_2+c_2/c_1$, so even the first compatibility equation is nontrivial.
A useful target is a normalization satisfying both finite-lattice deletion
and cyclotomic scale transport, with its comparison to Abel regularization
proved rather than stipulated.

\subsection{Analytic meaning of the formal Gamma-moment series}

The zero-radius result for $\KK$ rules out a naive analytic master germ, but
one-sided integrals can converge for suitable negative quadratic logarithmic
parameters. Study sectorial or distributional realizations whose asymptotic
coefficients are exactly~\eqref{eq:K}. Any such realization should preserve
the already proved fixed-content identities and must specify its contour
and parameter domain. No summability assertion is made here.

\subsection{Exact moment transforms and formal verification}

The beta-moment law diagonalizes the unweighted depth hierarchy under a
Laplace transform in $-\log(1-z)$. Investigate whether analogous transforms
exist for one or more spectral logarithmic decorations and whether they
close under a finite polygamma differential system. Independently, the
finite coefficient recurrence, the set-partition identity and the
Gamma-conjugation rule are suitable first targets for Lean formalization.
The analytic layer would then require formal proofs of the compensated
product and the polynomial Laplace--Stieltjes transfer. The numerical suites
in this package are not substitutes for those formal developments.

\section*{Conclusion}
\addcontentsline{toc}{section}{Conclusion}

A single normally convergent harmonic product, followed by a finite Gamma
operator, determines every cutoff and Abel polynomial in the stated
symmetrized logarithmic family. The formulas expose precisely when
positive-index Stieltjes constants survive and when they cancel. Diagonal
spectral restriction introduces additional terms, and root extraction
changes the endpoint scale; both effects have finite exact correction laws.
On the convergent side, beta-weighted moments and ordinary primitives admit
all-depth evaluations, while logarithmic primitives at integer shifts reduce
to a Gamma quotient with an explicit finite deletion. These are proved
identities with stated domains, not numerical reductions inferred from a
small residual.

\begin{thebibliography}{99}

\bibitem{DLMF}
NIST Digital Library of Mathematical Functions.
Sections 5.7 (Gamma series), 5.12 (beta function), 15.4 (hypergeometric
special values), 25.11 (Hurwitz zeta) and 25.14 (Lerch transcendent). Accessed 10 October 2026.
\url{https://dlmf.nist.gov/5.7};
\url{https://dlmf.nist.gov/5.12};
\url{https://dlmf.nist.gov/15.4};
\url{https://dlmf.nist.gov/25.11};
\url{https://dlmf.nist.gov/25.14}.

\bibitem{Hoffman}
M.~E.~Hoffman, Multiple harmonic series,
\emph{Pacific Journal of Mathematics} \textbf{152} (1992), 275--290.
The symmetric-sum attribution is also discussed in~\cite{KS}.

\bibitem{IKZ}
K.~Ihara, M.~Kaneko and D.~Zagier,
Derivation and double shuffle relations for multiple zeta values,
\emph{Compositio Mathematica} \textbf{142} (2006), 307--338.
Preprint record: \url{https://archive.mpim-bonn.mpg.de/id/eprint/725}.

\bibitem{KS}
M.~Kaneko and M.~Sakata, On multiple zeta values of extremal height,
arXiv:1505.01014. Consulted version contains the height-one Gamma generating
identity and the cutoff/shuffle regularization conventions.
\url{https://arxiv.org/abs/1505.01014}.

\bibitem{Aomoto}
K.~Aomoto, Special values of hyperlogarithms and linear difference schemes,
\emph{Illinois Journal of Mathematics} \textbf{34} (1990), 191--216.
Bibliographic attribution and the relevant generating identity are recorded
in~\cite{KS}.

\bibitem{Drinfeld}
V.~G.~Drinfeld, On quasitriangular quasi-Hopf algebras and a group closely
connected with $\operatorname{Gal}(\overline{\Q}/\Q)$,
\emph{Leningrad Mathematical Journal} \textbf{2} (1991), 829--860.
See~\cite{KS} for the height-one generating-function attribution.

\bibitem{repo}
ProveIt repository, \emph{Canonical polylogarithms manuscript},
with editorial records in
\src{Analysis/Polylogarithms/docs/manuscript/}, ProveIt repository.
Reviewed snapshot:
\src{fc4d3bf80534ad7c901d3b8c9e71baf2df0064ed}.
\url{https://github.com/VladimirReshetnikov/ProveIt/tree/fc4d3bf80534ad7c901d3b8c9e71baf2df0064ed/Analysis/Polylogarithms/docs/manuscript}.

\bibitem{repoalt}
ProveIt repository, \emph{Alternating harmonic polylogarithms},
report README and collection under
\src{Analysis/Polylogarithms/docs/reports/alternating-harmonic-polylogarithms/}.
README inspected at the same repository snapshot as~\cite{repo}.
\url{https://github.com/VladimirReshetnikov/ProveIt/tree/fc4d3bf80534ad7c901d3b8c9e71baf2df0064ed/Analysis/Polylogarithms/docs/reports/alternating-harmonic-polylogarithms}.

\bibitem{intake}
ProveIt repository, Incoming reports: intake and placement instructions,
\src{docs/incoming/README.md}, same reviewed snapshot.
\url{https://github.com/VladimirReshetnikov/ProveIt/blob/fc4d3bf80534ad7c901d3b8c9e71baf2df0064ed/docs/incoming/README.md}.

\end{thebibliography}
\end{document}
